\documentclass{article}
\usepackage{amsmath}
\usepackage{cancel}
\usepackage{graphicx} % for creating figures

\newcommand{\prob}{\mathcal{P}}


\title{Entropy}

\begin{document}
\maketitle
To begin, we define a \textbf{closed system} as system with constant energy,
constant number of particles, constant volume, and constant values for all
external parameters that affect the system (gravity, EM fields, etc...)

Quantum states are \textbf{accessible} when they are compatible with the
physical specification of the system. States that are not accessible are said to
have $0$ probability of being measured.

\subsubsection*{Probability}
Consider a closed system. Then we define the probability of finding the system
in a particular state $s$ to be given by
\begin{equation}
  \prob(s) = 1/g(s) 
\end{equation}
subject to the (unsurprising) condition that
\begin{equation}
  \sum_s \prob(s) = 1
\end{equation}

From this definition, we find that the average of observations for some
observable $X$ is given by
\begin{equation}
  \langle  X \rangle = \sum_s X(s)\prob(s) = \sum_s X(s)(1/g(s))
\end{equation}
Where the final equality is for a closed system. In taking this average, we
image a collection of similar systems with one in each accessible quantum state.
Such a collection is called \textbf{an ensemble}.


\subsubsection*{Most probable configuration}
Consider two separate systems $\mathcal{S}_1$ and $\mathcal{S}_2$ which are
brought into contact so that their energies are allowed to transfer freely
between the two systems. Such a scenario is called \textbf{thermal contact}. If
each system separately had energy $U_1$ and $U_2$, then the combined system is
closed with energy $U=U_1+U_2$.

\textbf{Question:} what determines how energy flows between the systems? 

Certainly, the solution is not so simple as \textit{the system with greater
  energy transfers to the system with less} because, intuitively, if we consider
two systems where the second system is just two copies of the first (i.e. we
have doubled U, N, and V) we expect nothing to happen. It must be more
complicated...

Instead, if we think about the number of accessible states, then the most
probably division of energy should occur when the combined system has the
greatest number of accessible states. To make this precise, let's think about
the multiplicity of the combined system. For every state specified when
$\mathcal{S}_1$ is at $U_1$, we expect there to be $g(N_2, U-U_1)=g(N_2,U_2)$
possible states for $\mathcal{S}_2$. Therefore, the combined multiplicity is
given as
\begin{equation}
  g(N,U) = \sum_{U_1}g(N_1, U_1)g(N_2, U-U_1)
\end{equation}
So that $g$ depends on $N_1, N_2,U_1, U_2 = U-U_1$. To simplify our lives, we    
model a system by the properties of it's most possible configuration, that is,
the choice of $U_1$ and $U_2$ for which $g(N,U)$ is maximized. Maxima occur when
the value of the function does not change for small changes in the arguments. In
other words, we must find the condition for which $dg=0$ when $dU_1$ $dU_2$
occur due to a small transfer of energy. The total differential of $g$ is given
by the multi-variable chain rule to be
\begin{align}
  dg &= \left( \frac{\partial g}{\partial U_1} \right)_{U_2}dU_1+\left( \frac{\partial g}{\partial U_2} \right)_{U_1}dU_2 \\
     &=  g_2\left( \frac{\partial g_1}{\partial U_1} \right)_{U_2}dU_1+g_1\left( \frac{\partial g_2}{\partial U_2} \right)_{U_1}dU_2\\
     &= 0\\
 \Rightarrow& \frac{1}{g_1}\left( \frac{\partial g_1}{\partial U_1} \right)_{U_2}dU_1  = -\frac{1}{g_2}\left( \frac{\partial g_2}{\partial U_2} \right)_{U_1}dU_2
\end{align}
But we know something about the relationship between $U_1$ and $U_2$! For a
closed system $U=U_1+U_2=\textrm{ constant}$. This means that $dU=dU_1+dU_2=0$
so that
\begin{equation}
  \frac{1}{g_1}\left( \frac{\partial g_1}{\partial U_1} \right)_{U_2}  = \frac{1}{g_2}\left( \frac{\partial g_2}{\partial U_2} \right)_{U_1}
\end{equation}
Now if we note that $1/f\dfrac{df}{dx}=\frac{d \ln(f(x))}{d x}$ by reverse
application of the chain rule, we have found
\begin{equation}
  \frac{\partial \ln g_1}{\partial U_1} = \frac{\partial \ln g_2}{\partial U_2}
\end{equation}

These quantities are equal for a system in thermal equilibrium so we should give
them a name. For historic reasons, we include a constant $k_B$, and define the
\textbf{Entropy} as
\begin{equation}
  S = k_b\ln g(N,U)
\end{equation}

The punchline of our derivation is \textbf{Two systems are in equilibrium when
 the slope of their S vs U graphs is the same.}













\end{document}
